\section{Propagation of regularity and the comparison error}\label{sec:regularity}

We now remove the outer regularity hypothesis from the short comparisons.
Throughout this section, $T>1$ and a sufficiently small $m>0$ are fixed,
and $\delta=m/10000$.  All limits are taken as $n\to\infty$ with these
parameters fixed.  Constants may depend on $T,m$, on fixed bounds for the
total covariance clocks, and on the finite recursion depth specified below.
No assertion in this section is uniform for a varying choice $m=m(n)$.

The two inputs are profile rigidity and the conditional comparison error.
Under the data
conditions \eqref{eq:3}--\eqref{eq:5}, an optimizer with an outer regular
region above a small scale $x$ satisfies the direct profile conclusion
of \Cref{prof:direct-profile}, provided
\[
 x\gtrsim a,\qquad e(x)x^2\gtrsim nE^2,
 \qquad e(x)\gtrsim x^2.
\]
At macroscopic scales the corresponding conclusion is
\Cref{prof:macro-profile}.  The possible slopes in the small
scale limit \eqref{eq:28} have fixed positive upper and lower bounds;
the macroscopic limiting quantile has the same kind of bounds relative
to mass.  These bounds depend only on $T$ and the uniform total clock
bounds.  Proposition~\ref{comp:conditional-error} bounds the upper
derivative per unit $\log\lambda$ in Lemma~\ref{comp:envelope} by
$n^{10\delta}E$ for
$\lambda\ge n^{-10}$ if the optimizer has an outer regular region above
\[
 A(\lambda)=n^\delta\frac{\sqrt n E}{\sqrt\lambda}.
\]
If $A(\lambda)$ is bounded below by a positive constant, that error
bound does not require a profile hypothesis.  The omitted initial
interval is allowed in each of these statements.

The propagation follows a hypothetical failure. At a selected bad
parent scale the parent has a regular outer region. Sufficient local
strength is excluded by a direct profile test. In the remaining case,
a successful child error estimate would give agreement with the
parent. When $x\to0$, this supplies the child's outer regularity for
the direct profile test; when $x$ stays bounded below, the macroscopic
direct result applies. In either case the child profile transfers
back and contradicts the parent failure. Thus the child error must
fail; that failure selects a strength and minimizer with a new bad
profile. The construction below raises $E$ by at least $n^{\delta/4}$
at each stage, while the data confine $E_j$ between fixed
stage-dependent multiples of $n^{-5/6+2m}$ and $n^{-1/2}$. A depth
fixed before $n\to\infty$ therefore suffices. The earlier finite-degree
recursion of Lemma~\ref{lem:cavity-recursion} has already been completed.

\subsection{Admissible families and child data}

\begin{definition}[Admissible comparison families]\label{reg:families}
The working mass interval is $I=[u_0,1]$, where either $u_0=0$ or
$u_0=o(a)$.  For a feasible quantile $r$ satisfying \eqref{eq:5}, write
\[
 P_r=\bigl(b_0-e(r)\mathbf 1_I,
                 h_0+e(r)r\mathbf 1_I\bigr),\qquad
 \mathcal P_e(r)=\frac14\int_I e(r)r^2\,du.
\]
An ordinary comparison family has an admissible bare endpoint
$(b_0,h_0)$ and admissible changed paths $P_r$ for every feasible $r$
and $0\le\lambda\le d^2$.  Its minimized value is
$\mathcal V(\lambda)=\min_r\{f_n(P_r)+\mathcal P_e(r)\}$.
A fixed strength family requires admissibility of every $P_r$ only
at the indicated positive strength.  Its bare coefficients need not
themselves define an admissible covariance path.  Adding a term
independent of $r$ does not change the minimizers or the profile
conclusions.  The notation $\mathcal V(\lambda)$ always denotes the
displayed value without such an added term; all error bounds and
strength derivatives refer to this value.  Thus a term that depends
on $\lambda$ must be subtracted at each strength before those bounds
are applied.
We require $d^2\le T-\rho$ for some fixed $\rho\in(0,T)$.
The child coefficient $c_0$ introduced below is always chosen with
$c_0^2\le T-\rho$, so this hypothesis persists in descendants.

The recursive comparisons below obey one of two covariance conditions.
Either they act on all masses with a fixed positive lower bound on
the bare matrix covariance, or they act only on $I$, with a fixed
positive matrix covariance on $I$ and a fixed positive matrix gap
at $u_0$.  A high-only family may be followed by all-mass families
when the already changed endpoint has a fixed positive matrix lower
bound.  Once that change is made, all later descendants act on all
masses.  High-only descendants have lower quantile endpoint zero.
Only the initial family may require a positive wall
$r_{\min}$ as in \eqref{eq:5}.
At the initial stage the selected covariance alternative must have
a fixed margin after the initial comparison as well.  For example,
an all-mass bare matrix lower bound $T/2$ and $d^2\le T/4$ give
such a margin.  In the high-only case the analogous requirement is
a fixed residual high covariance and high-to-low gap.  Alternatively,
a fixed global lower bound after the initial high-only comparison
allows the immediate change to all-mass descendants.
\end{definition}

The fixed positive bounds in this definition concern the initial
available covariance.  We will spend less than half of each such
bound along the entire finite chain.  Since $e$ is decreasing and
$re(r)$ is increasing, a comparison preserves monotonicity inside
its working interval.  The lower bound or the gap pays for matrix
nonnegativity and the matrix jump at its left endpoint.  A high-only
descendant adds zero field at that endpoint, so it preserves the
field boundary inequality already present in its bare endpoint.

\begin{lemma}[Data for a child comparison]\label{reg:child-data}
Suppose a parent comparison has data \eqref{eq:3}--\eqref{eq:5} and
an optimizing quantile with an outer regular region above $x$, where
\[
 A(\lambda)\le x\le1,\qquad e(x)=o(x^2).
\]
Fix $c_0>0$, independent of $n$, and set
\[
 d'=c_0x,\qquad
 E'=\min\bigl\{n^{-1/2}d'^2,
                      n^{-100\delta}e(x)x^3\bigr\},\qquad
 a'=\frac{\sqrt n E'}{d'}.
\]
Then, for all sufficiently large $n$, the child data satisfy
\eqref{eq:3}, with the same $m$, and
\[
 E'\ge n^{\delta/4}E,\qquad a'\ge a.
\]
The changed parent path, used as the child's bare endpoint, satisfies
\eqref{eq:4} at all its required test scales.  An omitted mass interval
of length $u_0=o(a)$ remains $o(a')$.
\end{lemma}

The two upper bounds defining $E'$ serve different purposes.  The
ceiling $n^{-1/2}d'^2$ keeps the child in the parameter range
\eqref{eq:3}.  The other bound makes a successful child error
negligible compared with the cost $e(x)x^3$ of a parent--child
separation.  The proof shows that both bounds still permit $E'$ to
be larger than $E$ by a fixed power of $n$.

\begin{proof}
The data conditions imply, with fixed implicit constants,
\[
 d\ge a\gtrsim n^{-1/6+m},\qquad
 E\gtrsim n^{-5/6+2m},\qquad nEa\gtrsim n^{3m}.
\]
The strength profile obeys $e(t)\gtrsim
\min(\lambda,d^3/t)$, while $te(t)$ is increasing.  Since
$x\ge A(\lambda)$ and $\sqrt\lambda\sqrt n E\le d^3$,
\[
 x\ge n^\delta a,\qquad
 e(x)x^2\gtrsim
 \min(\lambda x^2,d^3x)\gtrsim n^{1+\delta}E^2.
\]
These estimates are valid whether $x$ tends to zero or is macroscopic.

Put $E'_{\mathrm c}=n^{-1/2}c_0^2x^2$ and
$E'_{\mathrm s}=n^{-100\delta}e(x)x^3$.  For the lower bound in
\eqref{eq:3}, separately for the two candidate values, we have
\[
 \frac{E'_{\mathrm c}}{n^{-2/3+m}d'}
   =c_0 n^{1/6-m}x\gtrsim c_0n^\delta,
 \qquad
 \frac{E'_{\mathrm s}}{n^{-2/3+m}d'}
   \gtrsim c_0^{-1}n^{3m-99\delta}.
\]
Both tend to infinity; the required upper bound holds by definition.
The same calculation gives
\[
 \frac{E'_{\mathrm s}}E
 \gtrsim n^{1-99\delta}Ex\gtrsim n^{3m-99\delta}.
\]
Because $e(x)=o(x^2)$, the preceding lower bound on $e(x)x^2$
also gives $x^4/(n^{1+\delta}E^2)\to\infty$.  Consequently
\[
 \frac{E'_{\mathrm c}}E
     =c_0^2\frac{x^2}{\sqrt n E}
     \gg n^{\delta/2}.
\]
Taking the smaller candidate proves $E'/E\ge n^{\delta/4}$
eventually.  In the ceiling case $a'=c_0x\ge a$ eventually.
In the other case,
\[
 \frac{a'}a
 =c_0^{-1}n^{-100\delta}\frac{e(x)x^2d}E
 \gtrsim c_0^{-1}n^{1-99\delta}Ed
 \gtrsim n^{3m-99\delta},
\]
which proves the asserted cutoff inheritance.

It remains to check the actual covariance perturbations in
\eqref{eq:4}.  The parent outer bound and monotonicity give
$r(u)=O(y)$ when $u\asymp y\gtrsim x$, and $r(u)=O(x)$ when
$u=O(x)$.  Also $e(0)=o(x^2)$.  Indeed, along any subsequence on
which $d/x\to0$ this follows from $e(0)\le d^2$; if $d/x$ stays
bounded below, the argument $\lambda x/d^3\le x/d$ of $f$ stays
bounded, so $e(0)\asymp e(x)=o(x^2)$.  Thus at every scale
$y\gtrsim x$ the new matrix change is $o(y^2)$ and the new field
change is $o(y^3)$.

Consider now a test sequence $a'\lesssim y\ll x$.  This can only
occur in the nonceiling case.  Here
\[
 \frac{a'^3}{e(x)x}
 =c_0^{-3}n^{3/2-300\delta}e(x)^2x^5
 \gtrsim n^{7/2-300\delta}E^4x
 \gtrsim n^{9m-300\delta}.
\]
In particular the ratio has a polynomial divergence.  If
$\lambda x/d^3$ stays bounded, $e(0)\asymp e(x)$ and, for
$r\le Cx$,
\[
 e(r)\le O(e(x))=o(a'^2),\qquad re(r)\le O(e(x)x)=o(a'^3).
\]
The first conclusion uses also $a'\le d'=c_0x$.  Otherwise
$e(x)x\gtrsim d^3$, and the global profile bounds give
\[
 e(r)\le d^2=o(a'^2),\qquad
 re(r)\lesssim d^3\log n=o(a'^3).
\]
The polynomial divergence absorbs the logarithm.  These two
subsequence alternatives prove the required little-oh estimates at
every $y\gtrsim a'$.  The old bare endpoint already satisfies
\eqref{eq:4} there, since $a'\ge a$.  The omitted interval assertion
is immediate.  The same reasoning applies after a change from
high-only to all-mass comparisons: positive comparable test masses
at the new cutoff still lie above the original omitted interval.
\end{proof}

\subsection{Admissible movement and successful children}

To compare the parent minimizer with the child's, we move the parent
only part of the way toward the child.  The next lemma measures the
gain from any feasible such movement.  The barrier construction then
produces a movement without violating the quantile constraints.

\begin{lemma}[Movement using admissible paths]\label{reg:admissible-movement}
Let $r$ minimize the parent objective at a fixed positive strength,
and let a child comparison act on $J\supseteq I$.  Denote its final
profile by
\[
 \mathcal D(z)=d'^2 f(z/d'),
\]
and its final minimizing quantile by $z$.  Suppose its error bound is
successful, so the admissible doubly changed path
\[
 Q=\bigl(P_{r,b}-\mathcal D(z)\mathbf1_J,
              P_{r,h}+\mathcal D(z)z\mathbf1_J\bigr)
\]
satisfies
\[
 f_n(Q)+\frac14\int_J\mathcal D(z)z^2\,du
             \le f_n(P_r)+\mathrm{err}.
\]
Let $\widehat r$ be a feasible parent trial, equal to $r$ outside
the common working region, and suppose at every moved point
\[
 \widehat r=r+\omega(z-r),\qquad
 0\le\omega<\frac{\mathcal D(z)}{e(r)+\mathcal D(z)}.
\]
If its changed path $P_{\widehat r}$ is admissible, then
\[
 \int_I e(r)\omega(z-r)^2\,du\le4\,\mathrm{err}.
\]
This conclusion does not require an admissible bare parent endpoint
or a parent error bound.
\end{lemma}

\begin{proof}
Write $e=e(r)$, $\widehat e=e(\widehat r)$ and
$D=\mathcal D(z)$ in this proof only, and set
\[
 a=e\mathbf1_I+D\mathbf1_J-\widehat e\mathbf1_I,
 \qquad
 c=er\mathbf1_I+Dz\mathbf1_J-\widehat e\widehat r\mathbf1_I.
\]
At a moved point the log-slope bound $0\le-re'/e\le1$ implies
$\widehat e\le e/(1-\omega)$.  This follows from monotonicity
when $\widehat r\ge r$, and otherwise from
$\widehat r\ge(1-\omega)r$ and the monotonicity of $te(t)$.
Hence $a>0$ there.  At unmoved points of $J$, $a=D$ and $c=Dz$;
outside $J$, $a=c=0$.

The ratio needed below is integrable even if $a$ approaches zero.
Indeed put $A=e+D$, $\mu=D/A$, and $\Delta=z-r$ at a moved point.
The preceding bound gives
\[
 a\ge\frac{A(\mu-\omega)}{1-\omega},\qquad
 \frac ca=r+\left(\omega+\frac{A(\mu-\omega)}a\right)\Delta.
\]
The coefficient of $\Delta$ lies in $[\omega,1]$, so $c/a$ lies
between $\widehat r$ and $z$, both in $[0,1]$.  Thus
$c^2/a\le a\le e+D$.  At unmoved points of $J$ the ratio is
$Dz^2$, and outside $J$ it is zero.  These bounds also cover radius
zero and require no uniform lower bound on $a$.

The convex segment from $Q$ to $P_{\widehat r}$ consists of
admissible covariance paths.  Its pointwise variation is $(a,-c)$;
the variation itself need not be monotone.  By \eqref{eq:1} and
$B\ge S^2$, at every point of that segment,
\[
 -\frac a4B+\frac c2S\le\frac{c^2}{4a}.
\]
With the ratio interpreted as zero when $a=c=0$, integration gives
\[
 f_n(P_{\widehat r})\le f_n(Q)+\frac14\int\frac{c^2}a\,du.
\]
Combining this with parent minimality and the displayed child bound
proves
\[
 \int\left[er^2\mathbf1_I+Dz^2\mathbf1_J
       -\widehat e\widehat r^2\mathbf1_I-\frac{c^2}a\right]du
       \le4\,\mathrm{err}.
\]
Any fixed low-group penalty has canceled.  At every unmoved point
of $J$ the integrand is zero, including points below $u_0$ when
$J=[0,1]$.  At a moved point its value is
\[
 \left\{\frac{eD}{e+D}
  -\frac{\widehat e(e+D)}{e+D-\widehat e}
       \left(\omega-\frac D{e+D}\right)^2\right\}(r-z)^2.
\]
The expression in braces decreases as $\widehat e$ increases.
Substitution of its upper bound $e/(1-\omega)$ reduces it exactly
to $e\omega$, proving the assertion.  Only the three admissible
paths $P_r,Q,P_{\widehat r}$ have been evaluated.
\end{proof}

\begin{lemma}[Feasible barriers]\label{reg:barriers}
In the setting of Lemma~\ref{reg:child-data}, let $y\gtrsim x$.
Suppose, on an interval of mass length comparable to $y$ at masses
comparable to $y$, the parent and child quantiles are separated by
two fixed thresholds whose distance is comparable to $y$.
Assume $r\le Cy$ there, but impose no upper bound on $z/y$.
There is a feasible parent movement as in
Lemma~\ref{reg:admissible-movement} for which
\[
 \int e(r)\omega(z-r)^2\,du
       \ge c\min\{e(y),\mathcal D(y)\}y^3.
\]
Constants may depend on the separation thresholds and on the
fixed comparability constants.
\end{lemma}

\begin{proof}
In forward coordinates the constraints \eqref{eq:5} read
\[
 l(u)\le r'(u)\le\frac1{k(r(u))},\qquad l(u)=\frac1{M(u)}.
\]
Their forced lengths satisfy
\[
 \int_{u_0}^1l(u)\,du
    \le n^{-\delta}L\log\frac{1+u_*}{u_0+u_*}=o(x),
 \qquad
 \int_0^1 k(t)\,dt\le n^{-\delta}u_*=o(x).
\]
Also $l\le n^{-\delta}$, $k\le n^{-3\delta}$, and
$r_{\min}\le n^{-\delta}u_*=o(x)$.  The estimates follow directly
from \eqref{eq:5} and $u_*/a\lesssim n^{-3m}$.
We may shorten the separation interval by a fixed proportion to
leave room for ramps and strict threshold inequalities.  Put
\[
 t_y=\min\{1,\mathcal D(y)/e(y)\}.
\]

First suppose $z$ is above the upper threshold and $r$ below the
lower threshold.  Starting from $r_{\min}$, construct a barrier
$w$ with slope $l$ until an early portion of the separation
interval.  On that portion give it a fixed sufficiently large
slope to reach an intermediate threshold, and thereafter slope
$l$ again.  The ramp needs only a fixed small fraction of the
interval, since its rise is $O(y)$.  Its fixed slope is eventually
less than $n^{3\delta}$ and greater than $l$.  The $o(x)$ total
forced rise leaves $w<1$ and keeps a fixed distance of order $y$
between $w$ and $z$ wherever $w>r$.  Before the ramp the minimum
slope condition gives $w\le r$; after it, monotonicity of $z$
preserves the separation.  Outer regularity shows that clipping
can occur only at masses $O(y)$.

Let $v=\max(r,w)$ and set
$\widehat r=(1-t)r+tv$, where $t=\eta t_y$ and $\eta>0$ is a
sufficiently small fixed constant.  Both endpoints of the quantile
are unchanged.  The lower slope bound is preserved.  Where clipping
occurs, $v'=w'\le n^{3\delta}\le1/k(r)$; hence
\[
 \widehat r'\le\frac1{k(r)}\le\frac1{k(\widehat r)}.
\]
Thus the upper bound is preserved as well.

We verify the movement fraction even if $z/y$ is unbounded.
At clipped points $r<w\le C_1y$ and $z-w\ge c_1y$, so
$z-r\ge c_2z$.  Monotonicity of $t\mathcal D(t)$ and the
fixed-scale bounds for $\mathcal D$ give
\[
 \mathcal D(z)(z-r)\ge c_3y\mathcal D(y),
 \qquad \mathcal D(z)\le C_2\mathcal D(y).
\]
If $y/x\to\infty$, outer regularity gives $r\asymp y$ at the
clipped masses, and $e(r)\le C_3e(y)$.  Therefore the largest
allowed absolute displacement is bounded below by
\[
 \frac{\mathcal D(z)(z-r)}{\mathcal D(z)+e(r)}
       \ge c_4yt_y.
\]
If instead $y\asymp x$, then
$e(r)\le e(0)=o(x^2)$ and $\mathcal D(y)\asymp x^2$.
The same bound holds, with $t_y=1$ eventually, without a lower
bound on $r/y$.  Every sequence has a subsequence in one of
these two cases.  Choosing $\eta$ small makes the movement
strictly admissible.  On a remaining interval of length $cy$,
both $v-r$ and $z-r$ are at least $cy$, and $e(r)\ge ce(y)$.
Since $\omega(z-r)=\widehat r-r=t(v-r)$, the gain there is at
least $c e(y)t_y y^3$.

For a downward separation, start an upper barrier $w$ at an
intermediate positive threshold plus the accumulated slope $l$.
Keep that slope until a late portion of the separation interval.
There let $w'=1/k(w)$ until $w$ meets
\[
 U(u)=1-\int_u^1l(v)\,dv,
\]
and follow $U$ thereafter.  Every feasible $r$ satisfies $r\le U$.
The ramp takes $o(x)$ mass: while $w<U$, the speed of $U$ is at
most $n^{-\delta}$ and that of $w$ is at least $n^{3\delta}$,
so its length is bounded by twice $\int_0^1k=o(x)$ for large $n$.
Place this entire ramp inside the reserved late portion of the
separation interval; its length is $o(x)=o(y)$, so it fits there.
Where $v=\min(r,w)$ differs from $r$, both $v$ and $r$ lie in
$[cy,Cy]$, while $z$ is a fixed distance below $v$.  The assertion
before the separation interval follows from monotonicity of $z$;
the barrier is no longer active after its ramp.  It agrees with
$r$ at the two endpoints, because $r_{\min}=o(x)$ and $r\le U$.

Set $\psi(q)=\int_0^qk(t)\,dt$ and now interpolate by
\[
 \psi(\widehat r)=(1-t)\psi(r)+t\psi(v),\qquad t=\eta t_y.
\]
Both $r$ and $v$ satisfy the two slope bounds.  Therefore
\[
 k(\widehat r)\widehat r'
      =(1-t)k(r)r'+tk(v)v'\le1.
\]
For the other bound write $k(q)=A_0/(q+B_0)^2$, with
$A_0=n^\delta u_*^2$ and $B_0=n^{2\delta}u_*$.  Then
\[
 k\bigl(\psi^{-1}(s)\bigr)
       =\frac1{A_0}\left(\frac{A_0}{B_0}-s\right)^2
\]
is convex.  Jensen's inequality gives
\[
 k(\widehat r)\widehat r'
    \ge l\{(1-t)k(r)+tk(v)\}\ge l k(\widehat r).
\]
This proves feasibility.  On the clipping region the radii are
comparable to $y$, so the distortion between the interpolation
fraction in $\psi$ and the fraction in radius is bounded above
and below by fixed constants.  Here $e(r)\asymp e(y)$ and
$\mathcal D(z)\gtrsim\mathcal D(y)$, since $z\le Cy$.
Taking $\eta$ smaller if necessary gives the required strict
movement fraction.  On an interval of length $cy$, the radius
displacement is at least $ct_yy$, and $r-z\ge cy$.
The same gain bound follows.

Both constructions preserve the parent's lower wall.  Hence their
changed paths are admissible in every family of
Definition~\ref{reg:families}.  All moved masses lie above both
lower working boundaries, since $u_0=o(a)\le o(x)$.
\end{proof}

\begin{lemma}[A successful child error excludes a parent failure]
\label{reg:successful-child}
In Lemma~\ref{reg:child-data}, suppose the child's final error is at
most $n^{20\delta}E'$.  If $x\to0$, every sequence has a subsequence
on which the parent quantile satisfies \eqref{eq:28}, with the bounds
on $\kappa$ stated there.  If $x$ stays bounded below by a positive constant,
it converges locally uniformly on $(0,1)$ to the Parisi quantile $q_T$.
\end{lemma}

\begin{proof}
First, on every test scale $y\gtrsim x$, the parent and final child
quantiles agree to $o(y)$ in normalized mass measure on compact
positive mass bands of scale $y$.  At macroscopic scales take bands
inside $(0,1)$.  If this failed, monotone subsequence compactness
(truncating $z/y$ first if necessary) would produce an interval of
length comparable to $y$ on which the quantiles are separated by
fixed thresholds of distance comparable to $y$.  The parent outer
bound and monotonicity give $r\le Cy$ there.  Thus
Lemmas~\ref{reg:admissible-movement} and~\ref{reg:barriers} imply
\[
 c\min\{e(y),\mathcal D(y)\}y^3\le4n^{20\delta}E'.
\]
Both radius times strength profiles are increasing.  Moreover
$\mathcal D(x)=c_0^2f(1/c_0)x^2\asymp x^2$ and $e(x)=o(x^2)$.
Consequently, for every $y\gtrsim x$,
\[
 \min\{e(y),\mathcal D(y)\}y^3
       \gtrsim e(x)x^3,
 \qquad
 n^{20\delta}E'\le n^{-80\delta}e(x)x^3,
\]
a contradiction.  For $y/x$ bounded below by a fixed number
smaller than one, the same inequalities follow from the fixed-scale
profile bounds; thus the conclusion includes any fixed positive
multiple of $x$.

The child's numerical direct-test conditions now hold at $x$:
\[
 a'\le c_0x,\qquad
 nE'^2\le c_0^4x^4\lesssim\mathcal D(x)x^2,
 \qquad \mathcal D(x)\asymp x^2.
\]
Its changed paths and feasible competitors are admissible, and
$d'^2\le T-\rho$ by Definition~\ref{reg:families}.

Suppose $x\to0$.  This agreement gives the child a regular outer
region above $x$, with fixed, possibly enlarged constants.  Indeed,
a violation at masses whose ratio to $x$ tends to infinity gives,
by monotonicity and a neighboring interval of comparable masses,
a threshold separation from the parent's regular outer profile.
A violation at bounded multiples is included by enlarging the
lower cutoff of that outer region.
\Cref{prof:direct-profile} therefore applies to the
child, giving \eqref{eq:28}.  Agreement in measure and
monotonicity transfer this limit to the parent at every positive
scaled mass: bracket a point by nearby continuity points of the
strictly increasing linear limit and then shrink the bracket.
If $x$ stays bounded below, apply
\Cref{prof:macro-profile} to the child instead and use
the same agreement.  Near mass one, fixed loose comparison bounds
follow from monotonicity and $r\le1$.  This proves the assertion.
\end{proof}

\subsection{Finite propagation and the field wall}

\begin{proposition}[Finite propagation]\label{reg:finite-propagation}
For every sequence of comparison families in Definition~\ref{reg:families}
with data \eqref{eq:3}--\eqref{eq:5}, there are fixed positive
constants $c<C$ such that every minimizing quantile at a positive
strength $0<\lambda\le d^2$ satisfies
\[
 cu\le r(u)\le Cu\qquad
 \left(n^\delta\frac{\sqrt n E}{\sqrt\lambda}\le u\le1\right)
\]
for all sufficiently large $n$ along that sequence.  For ordinary admissible bare
endpoints, the entire comparison satisfies
\[
 0\le\mathcal V(\lambda)-f_n(b_0,h_0)
                  \le n^{20\delta}E,
       \qquad 0\le\lambda\le d^2.
\]
For a fixed strength family only the profile assertion is made;
its proof uses error estimates exclusively for admissible child
endpoints.
\end{proposition}

\begin{proof}
\emph{Fix the bounds and recursion depth.}
Enlarge the initial total-clock bounds by one to allow for descendant
field additions.  Use these enlarged bounds when choosing the
constants in the limiting profiles.
Choose $c,C$ strictly outside the bounds of all the limiting
profiles described at the start of this section, with additional
room near mass one.  Fix
\[
 N=\left\lceil\frac{10}{\delta}\right\rceil.
\]
Before taking $n\to\infty$, choose $c_0>0$ small enough that
$Nc_0^2<1$ and $Nc_0^2$ is less than half each initial fixed residual
matrix lower bound or gap.  It may also be decreased to meet the fixed bounds needed
in the direct tests, including $c_0^2\le T-\rho$.
Every child has $d'=c_0x\le c_0$, and hence $d'^2\le T-\rho$.
The sum of all further matrix removals along a chain of length
$N$ is therefore at most $Nc_0^2$.  All field additions are
increasing and have total size at most $Nc_0^2$.  Thus the total
clock bounds and the covariance alternatives of
Definition~\ref{reg:families} persist along this chain.  In
particular every child endpoint, every feasible child path, and
every segment used in Lemma~\ref{reg:admissible-movement} are
admissible.

\emph{A failed parent profile forces a failed child error.}
Assume first that the profile assertion fails along a sequence.
At its selected strength choose an offending mass $x\ge A(\lambda)$
within a factor two of the supremum of all offending masses.
There are no offending masses above $2x$, so the path has a regular
outer region above $x$.  The numerical direct-test conditions
$x\gtrsim a$ and $e(x)x^2\gtrsim nE^2$ follow from
\eqref{eq:8}.  If
$e(x)/x^2$ has a positive lower bound along a subsequence, Theorem~\ref{prof:direct-profile} or
Proposition~\ref{prof:macro-profile}, according to the scale, immediately
contradicts the choice of the loose constants $c,C$.
After subselection the only remaining case is $e(x)/x^2\to0$.
Lemma~\ref{reg:child-data} constructs a child from the changed
parent endpoint.  If its error estimate were successful,
Lemma~\ref{reg:successful-child} would again contradict the parent
violation.  Hence a profile failure forces a child error failure.

\emph{A failed error selects the next profile failure.}
Every such child has an admissible bare endpoint.  If all its
minimizers at all relevant strengths had the asserted profile,
Proposition~\ref{comp:conditional-error} and the optimized envelope
of Lemma~\ref{comp:envelope} would give
\[
 \mathcal V(d'^2)-f_n(b'_0,h'_0)
       \le O(\log n)n^{10\delta}E'+O(n^{-10})
       <n^{20\delta}E'
\]
eventually.  Strengths below $n^{-10}$ contribute $O(n^{-10})$
by bounded overlaps; when $A(\lambda)$ is macroscopic the
conditional error bound is unconditional.  The derivative bound
uses an old minimizer for positive increments, or equivalently
the Lipschitz envelope inequality, and needs no monotone selection
of minimizers.  It follows that a child error failure selects
a strength and a minimizing quantile with a profile failure.
The same argument applies if the initial failure was an error
failure at an ordinary endpoint.  The lower bound in the
proposition is the nonnegative segment interpolation in
\eqref{eq:6}.

\emph{The error parameter cannot grow for the whole fixed chain.}
We may therefore repeat the construction: each bad profile forces
a bad error for a valid child; that bad error selects the next bad
profile.  Lemma~\ref{reg:child-data} preserves all asymptotic
endpoint hypotheses and multiplies $E$ by at least
$n^{\delta/4}$ at each stage.  All choices of constants have
already been made.  At most $N$ successive subsequence extractions
are involved, so the eventual inequalities can be imposed together.
On the other hand \eqref{eq:3}, at every stage, confines the error
parameter to
\[
 c_j n^{-5/6+2m}\le E_j\le C_j n^{-1/2},
\]
where $c_j,C_j$ are fixed constants for that stage.  A chain of
$N$ increases would imply
$E_N\ge n^{N\delta/4}E_0$, contradicting these bounds since
$N\delta/4\ge5/2>1/3-2m$.  This proves both assertions.

For a fixed strength root, the argument begins directly with its
bad profile.  There is no root sweep.  Its first child starts from
the admissible changed path $P_r$, and every subsequent error
sweep has an admissible bare endpoint.  Lemma~\ref{reg:admissible-movement}
shows explicitly that no bare root pressure is needed.
\end{proof}

\begin{corollary}[Uniformity for the field wall]\label{reg:wall-uniformity}
Let the parent have data \eqref{eq:3}--\eqref{eq:5} and work on
$[s,1]$, with $s=o(a)$, at final strength
$e(q)=d^2f(q/d)$ and lower wall
$0<r_{\min}\le u_*/M(s)$.  Suppose $d=o(1)$ and $d^2=o(\varepsilon)$.
For $0<\varepsilon\le T$ and $0\le\theta\le1$, put
\[
 c_L=\frac12\min(\varepsilon,T-\varepsilon),\qquad
 r_0=\frac{d^2r_{\min}}{2\varepsilon},
\]
and take increasing $0\le\varphi\le1$ on $[0,s]$.  The bare
algebraic coefficients are
\[
 b_0=T-(\varepsilon+\theta c_L)\mathbf1_{u<s},\qquad
 h_0=\theta c_Lr_0\varphi(u)\mathbf1_{u<s}.
\]
The profile conclusion of Proposition~\ref{reg:finite-propagation}
holds uniformly along arbitrary sequences of these wall parameters,
including $\theta\to0$ and $T-\varepsilon\to0$.
\end{corollary}

\begin{proof}
The low matrix covariance is nonnegative because
$c_L\le(T-\varepsilon)/2$.  At the high boundary,
$e(r_{\min})=d^2$ for large $n$, since $r_{\min}/d\to0$.
The low field is at most
\[
 c_Lr_0\le\frac14d^2r_{\min},
\]
whereas the high field is at least $d^2r_{\min}$.
The high matrix gap is at least $\varepsilon-d^2>0$.
Thus every changed parent trial is an admissible path.
The low penalty, if included, is independent of the high quantile
and cancels in Lemma~\ref{reg:admissible-movement}.

If $\varepsilon<T/4$, the low matrix covariance is greater than
$5T/8$; for large $n$ the high covariance has a fixed positive
lower bound as well.  We therefore use all-mass descendants.
If $\varepsilon\ge T/4$, the high-to-low gap is at least
$T/4-d^2$, so we use high-only descendants with zero wall.
Choose the recursion budget as a fixed small fraction of $T$.
These alternatives remain valid even when the low covariance
tends to zero, since in that case descendants leave it unchanged.
Above $s$ the original bare coefficients are exactly $T,0$;
\eqref{eq:4} is therefore uniform there.  Lemma~\ref{reg:child-data}
preserves it in every descendant, and $s=o(a')$ persists.

Any failure of uniformity would yield a sequence of wall parameters
and, after subselection, one of these two fixed-margin alternatives.
Proposition~\ref{reg:finite-propagation} excludes that sequence.
Only changed paths are used, even when the displayed bare field
has a downward jump at $s$.
\end{proof}
